
\section{Proposal}
\label{Pro}

\begin{figure}[t]
  \centering
  \scalebox{1.0} {
    \Qcircuit @C=1em @R=1em {
  \lstick{\ket{x_1}} & \multigate{3}{F} & \qw           &          & \qw           & \multigate{4}{F \cdot a} & \qw             \\
  \lstick{\ket{x_2}} & \ghost{F}        & \qw           &          & \qw           & \ghost{F \cdot a}        & \qw             \\
                     &                  & \push{\vdots} & \push{=} & \push{\vdots} &                          & \push{\vdots}   \\
  \lstick{\ket{x_n}} & \ghost{F}        & \qw           &          & \qw           & \ghost{F \cdot a}        & \qw             \\
  \lstick{\ket{y}}   & \targ\qwx[-1]\qw & \qw           &          & \gate{H}      & \ghost{F \cdot a}        & \gate{H}       
}

  }
  \caption{Synthesize a single target gate from a controlled gate}
  \label{fig-tof}
\end{figure}

\subsection{Direct Synthesis of Phase Gates}
\label{phase-direct}

In a parallel to~\figref{fig-phase-n4}, Controlled-NOT gates can be implemented by
synthesizing the phase oracle for the function $F \cdot a$ where $F$ is the function
controlling the switching of the NOT, and $a$ is the value of an ancilla to be used as
the target. H-gates are then placed on either side of the inputs to the phase gate that
take in $a$ and output $a$. This is demonstrated in~\figref{fig-tof}. Controlled-NOT
functions such as the Toffoli gate are often constructed in the same way (A Toffoli can
be constructed by applying the same procedure to a three variable Controlled-Z gate).

In this manner, it is clear that to implement a Controlled-NOT gate controlled by a function
$F$ of $n$ variables, it is necessary to compute a phase oracle of $F\cdot a$ which has
$n+1$ variables. Generating such a circuit can come with an increase in T-count. For example
implementing $x_a \cdot x_b \oplus z$, which contains a 3-variable phase gate inside, costs
7 T-gates, while implementing the phase oracle for $x_a \cdot x_b$ costs no T-gates. In this
way, a savings in T-count can be realized merely by implementing the phase oracle
directly.

\subsection{CNOT+T Circuits and Boolean Fourier Transform}
\label{Mot:CNOTBool}

As already noted in \secref{Pre:Four}, the phase polynomial with the optimal T-count can
already be calculated. In many cases, this T-count is actually 0. Observe that
\figref{fig-series-all} can be implemented by putting an oracle for
$x_a \cdot x_b \cdot \bar{x_c}$ and $\bar{x_b} \cdot x_c \cdot x_d$ in series with each
other, and then expanding each of them out to their CNOT+T implementation. Doing so will
enable the sum of paths method to combine T-gates, such that the phases needed to implement
that term can combine into a value that does not need any T-gates to implement. Utilizing a
library of such functions where the number of variables $n \leq 3$, we can enable lower
T-counts through cancellations and by guaranteeing the lowest T-counts for each of the
individual components.

\subsection{Optimizing ESOP Expressions for T-count and T-depth Reduction}
\label{Mot:Lib}

We can deduce from~\secref{Pre:Four} that any cube of 3 literals or fewer can be synthesized
purely as a quantum circuit. Moreover, any cube of 2 literals or fewer can be realized
without using any T-gates. Knowing that, it is thus useful to rewrite any Boolean functions
with as few cubes of over four literals as possible. This allows us to reduce the number
of T-gates, as well as implement as much of the circuit as possible using only CNOT+T, which
allows the usage of tools such as sum of paths simplification in Sec.~\ref{Mot:CnotOpt} and
matroid partitioning from~\cite{bib-amy-matroid} to optimize for T-count and T-depth.

As an additional benefit, many ESOP minimization algorithms such as EXORCISM-4~\cite{bib-exorcism} minimizes the
sharing of variables between any two cubes. That means that pairing ESOP-based phase oracle generation
with ESOP minimization generates many blocks that can be executed in parallel, reducing the T-depth.

\begin{example}
  \label{ex-esop-synth2}
  Let's revisit Fig.~\ref{fig-esop-synth}. Observe that
  $\bar{x_a}\bar{x_c} \oplus \bar{x_a} x_b \bar{x_c}\bar{x_d} \oplus x_a x_b \oplus x_a \bar{x_c} x_d$
  can be rewritten as
  $\bar{x_a} \bar{x_b} \bar{x_c} \bar{x_d} \oplus \bar{x_c} x_d \oplus x_a x_b$.
  Since 2 literal cubes have a T-count and T-depth of 0, 3 literal cubes have a T-count of
  7 and T-depth of 3, and $n \geq 4$ can use~\figref{fig-phase-n4} to implement a cube by
  using a Multiple Control Toffoli (MCT) gate in place of the $F$-controlled gate. When
  implemented with a single ancilla, this circuit T-count of 15 and T-depth of 7. This
  means that after the rewrite, the function saves a further 7 T-gates, and a T-depth of 3.
\end{example}

\begin{figure}[t]
  \centering
  \scalebox{1.0} {
    \input{img/flowchart.tex}
  }
  \caption{The overall flow of the proposed method}
  \label{fig-flow}
\end{figure}

\begin{figure*}
  \begin{minipage}[t]{0.45\textwidth}
    \vspace{0pt}
    % Requires:
    % \usepackage{amsmath}
    % \usepackage{tikz}
    % \usetikzlibrary{positioning}
    \scalebox{0.8} {
      \begin{tikzpicture}[
          font=\small,
          box/.style={draw, rounded corners=2pt, align=left, inner sep=6pt},
          >=Latex
        ]
        \node[box] (og) {$\displaystyle
          \mathbf{OG:}\;
          \Bigl(x_{3} \bar{x_{1}} \bar{x_{2}} x_{6} \bar{x_{5}} \bar{x_{4}}\Bigr)
          \oplus \Bigl(x_{3} x_{6}\Bigr)
          \oplus \Bigl(x_{1} x_{5}\Bigr)
          \oplus \Bigl(x_{2} \bar{x_{1}} \bar{x_{6}} \bar{x_{4}} \bar{x_{3}} \bar{x_{5}}\Bigr)$\\[-1pt]
          $\displaystyle
          \oplus \Bigl(\bar{x_{5}} \bar{x_{2}} \bar{x_{6}} x_{1} \bar{x_{4}} x_{3}\Bigr)
          $};

        \node[box, below=12mm of og] (red) {$\displaystyle
          \mathbf{Reduced:}\;
          \Bigl(x_{1} x_{5}\Bigr)
          \oplus \Bigl(\bar{x_{1}} \bar{x_{2}} x_{3} \bar{x_{4}}\Bigr)
          \oplus \Bigl(\bar{x_{1}} \bar{x_{2}} \bar{x_{3}} x_{4} \bar{x_{5}}\Bigr)
          \oplus \Bigl(x_{3}\Bigr)
          $};
%%         \node[box] (og) {$\displaystyle
%%           \mathbf{OG:}\;
%%           \Bigl(x_{3} ! x_{1} ! x_{2} x_{6} ! x_{5} ! x_{4}\Bigr)
%%           \oplus \Bigl(x_{3} x_{6}\Bigr)
%%           \oplus \Bigl(x_{1} x_{5}\Bigr)
%%           \oplus \Bigl(x_{2} ! x_{1} ! x_{6} ! x_{4} ! x_{3} ! x_{5}\Bigr)$\\[-1pt]
%%           $\displaystyle
%%           \oplus \Bigl(! x_{5} ! x_{2} ! x_{6} x_{1} ! x_{4} x_{3}\Bigr)
%%           $};

%%         \node[box, below=12mm of og] (red) {$\displaystyle
%%           \mathbf{Reduced:}\;
%%           \Bigl(x_{1} x_{5}\Bigr)
%%           \oplus \Bigl(! x_{1} ! x_{2} x_{3} ! x_{4}\Bigr)
%%           \oplus \Bigl(! x_{1} ! x_{2} ! x_{3} x_{4} ! x_{5}\Bigr)
%%           \oplus \Bigl(x_{3}\Bigr)
%%           $};

        \draw[->, line width=0.8pt] (og.south) -- (red.north);

      \end{tikzpicture}
    }
    \subcaption{ESOP read-in and simplification}
    \label{fig-esop-simp}
  \end{minipage}
  \begin{minipage}[t]{0.45\linewidth}
    \vspace{0pt}
    \hspace{-1cm}
    \centering
    \scalebox{0.8} {
      % TikZ diagram for the given BLIF
% - Oriented horizontally: inputs on the left, outputs on the right
% - Each .names block is shown as a small LUT/AND/INV-style logic box
% - Wires follow the BLIF net dependencies

\begin{tikzpicture}[
  x=1.35cm, y=0.9cm,
  >=Latex,
  wire/.style={-Latex, line width=0.5pt},
  nodebox/.style={draw, rounded corners=2pt, align=center, inner sep=3pt, minimum height=8mm},
  port/.style={draw, circle, inner sep=1.2pt},
  lab/.style={font=\small}
]

% Inputs (left)
\node[lab] (x0) at (-1,  3) {$x_0$};
\node[lab] (x1) at (-1,  1.5) {$x_1$};
\node[lab] (x2) at (-1,  0) {$x_2$};
\node[lab] (x3) at (-1, -1.5) {$x_3$};
\node[lab] (x4) at (-1, -3) {$x_4$};

% Internal logic nodes (placed left-to-right following dataflow)

% .names x0 x1 new_n8   (truth table: 00 1)  => new_n8 = ~x0 & ~x1
\node[nodebox] (n8) at (1, 3) {\scriptsize $i_0 = !x_0  !x_1$};
\node[nodebox] (n8b) at (1, 1.5) {\scriptsize $i_1 = !x_0  !x_1$};
% .names x2 new_n8 new_n9  (11 1) => new_n9 = x2 & new_n8
\node[nodebox] (n9) at (2.5, 1) {\scriptsize $i_2 = x_2  i_0$};

% .names x3 new_n9 po0  (01 1) => po0 = ~x3 & new_n9
\node[nodebox] (po0box) at (4, 0) {\scriptsize $po_0 = !x_3  i_2$};

% .names x2 new_n8 new_n11 (01 1) => new_n11 = ~x2 & new_n8
\node[nodebox] (n11) at (2.5, -1) {\scriptsize $ i_3 = !x_2  i_1$};

% .names x3 new_n11 new_n12 (11 1) => new_n12 = x3 & new_n11
\node[nodebox] (n12) at (4, -2) {\scriptsize $i_4 = x_2  i_3$};

% .names x4 new_n12 po1 (01 1) => po1 = ~x4 & new_n12
\node[nodebox] (po1box) at (5.5, -3) {\scriptsize $po_1 = !x_4  i_4$};

% Wires from inputs to logic
\draw[wire] (x0.east) -- (n8.west);
\draw[wire] (x1.east) -- (n8.west);

\draw[wire] (x0.east) -- (n8b.west);
\draw[wire] (x1.east) -- (n8b.west);

\draw[wire] (x2.east) -- (n9.west);
\draw[wire] (n8.east) -- (n9.west);

\draw[wire, bend left=25] (x3.east) to (po0box.west);
\draw[wire] (n9.east)  -- (po0box.west);

\draw[wire] (x2.east) -- (n11.west);
\draw[wire] (n8b.east)  -- (n11.west);

\draw[wire] (x3.east) -- (n12.west);
\draw[wire] (n11.east) -- (n12.west);

\draw[wire] (x4.east) -- (po1box.west);
\draw[wire] (n12.east) -- (po1box.west);

% Output labels on the far right (ports)
\node[lab] ($po_0$) at (10, 0.75) {};
\node[lab] ($po_1$) at (10, -0.75) {};

% Draw explicit output wires to the right (so outputs are clearly on the right edge)
\draw[wire] (po0box.east) -- ++(1.2,0) node[lab, right] {po0};
\draw[wire] (po1box.east) -- ++(1.2,0) node[lab, right] {po1};

\end{tikzpicture}

    }
    \subcaption{$n \geq 4$ Subject Graph}
    \label{fig-ex-network-og}
  \end{minipage}
  \begin{minipage}[t]{0.45\linewidth}
    \vspace{0pt}
    \centering
    \scalebox{0.8} {
      \input{img/fig-ex-network}
    }
    \subcaption{Simplified Logic Network}
    \label{fig-ex-network}
  \end{minipage}
  \begin{minipage}[t]{0.45\textwidth}
    \vspace{0pt}
    \centering
    \scalebox{0.8} {
      \Qcircuit @C=0.5em @R=0.2em @!R { \\
                         &      &           &           & \push{n \geq 4} &           &            &          &           &          &     & \push{n \leq 3} &            &          &       \\
                         &      &           &           &                 &           &            &          &           &          &     &                 &            &          &       \\
  \lstick{\ket{{x}_{0}}} & \qw  & \ctrlo{5} & \qw       & \qw             & \qw       & \qw        & \qw      & \qw       & \qw      & \qw & \qw             & \ctrl{5}   & \qw      & \qw   \\
  \lstick{\ket{{x}_{1}}} & \qw  & \ctrlo{6} & \qw       & \qw             & \qw       & \qw        & \qw      & \qw       & \qw      & \qw & \qw             & \qw        & \qw      & \qw   \\
  \lstick{\ket{{x}_{2}}} & \qw  & \qw       & \ctrl{6}  & \qw             & \ctrl{5}  & \ctrlo{5}  & \qw      & \qw       & \qw      & \qw & \ctrlo{5}       & \qw        & \qw      & \qw   \\
  \lstick{\ket{{x}_{3}}} & \qw  & \qw       & \qw       & \ctrlo{5}       & \qw       & \qw        & \ctrl{5} & \qw       & \ctrl{5} & \qw & \qw             & \qw        & \gate{Z} & \qw   \\
  \lstick{\ket{{x}_{4}}} & \qw  & \qw       & \qw       & \qw             & \qw       & \qw        & \qw      & \ctrlo{5} & \qw      & \qw & \qw             & \qw        & \qw      & \qw   \\
  \lstick{\ket{{x}_{5}}} & \qw  & \qw       & \qw       & \qw             & \qw       & \qw        & \qw      & \qw       & \qw      & \qw & \qw             & \ctrl{-5}  & \qw      & \qw   \\
  \lstick{\ket{{x}_{6}}} & \qw  & \qw       & \qw       & \qw             & \qw       & \qw        & \qw      & \qw       & \qw      & \qw & \qw             & \qw        & \qw      & \qw   \\
  \lstick{\ket{0}     }  & \qw  & \targ\qw  & \ctrl{1}  & \qw             & \ctrl{1}  & \ctrl{1}   & \qw      & \qw       & \qw      & \qw & \ctrl{1}        & \qw        & \qw      & \qw   \\
  \lstick{\ket{0}     }  & \qw  & \qw       & \targ     & \ctrl{-3}       & \targ     & \targ      & \ctrl{1} & \qw       & \ctrl{1} & \qw & \targ           & \qw        & \qw      & \qw  \\
  \lstick{\ket{0}     }  & \qw  & \qw       & \qw       & \qw             & \qw       & \qw        & \targ    & \ctrl{-2} & \targ    & \qw & \qw             & \qw        & \qw      & \qw  \gategroup{3}{2}{14}{12}{.7em}{--} \gategroup{3}{13}{14}{15}{.7em}{.}\\
                         &      &           &           &                 &           &            &          &           &          &     &                 &            &          &       
}

    }
    \subcaption{Synthesize gates}
    \label{fig-ex-qc}
  \end{minipage}
  \caption{Synthesis Example}
  \label{fig-ex-pro}
\end{figure*}

\subsection{Proposed Method Overview}
\label{Subsec:ProMet}
In this section, we present the overall flow of our proposal, depicted in Fig.~\ref{fig-flow}. The blue boxes represent precomputation, while tan boxes represent runtime.

\begin{itemize}
\item \textbf{(Generate 3bit phase polynomials)} : Precompute library of phase polynomials of Boolean functions of
  3 bits or fewer using the Boolean Fourier transform.
\item \textbf{(Synthesize 3bit phase subcircuits)} : The phase polynomials are then synthesized as subcircuits
  using the methods introduced in~\secref{Pre:CnotPhase}. The results are stored in hard disk, to be loaded into
  memory at runtime.
\item \textbf{(Decompose $f$ into ESOP expression)} : At runtime, a Boolean function is read in and converted to
  an ESOP expression.
\item \textbf{(Minimize the number of $n \geq 4$ cubes in ESOP)} : The ESOP is optimized by adapting
  EXORCISM-4~\cite{bib-exorcism} to reduce the number of $n \geq 4$ cubes. This method is explained
  in Sec.~\ref{Pro:Minimize}.
\item \textbf{(Map the $n \leq 3$ cubes to the 3bit library)} : Takes all cubes of $n \leq 3$ and maps them to 
  appropriate functions from the 3 bit library, inserting the corresponding subcircuit into the phase oracle,
  as described in Sec.~\ref{Pro:n4}.
\item \textbf{(Synthesize logic network and QC from $n \geq 4$ cubes)} : Generates a logic network from all cubes
  of $n \geq 4$ as a logic network with the values of each of the cubes as separate outputs, and synthesizes the
  quantum circuit as a network of Toffoli gates and Controlled-Z gates, as described in Sec.~\ref{Pro:n4}.
\item \textbf{(Append $n \geq 4$ and $n \leq 3$)} : The results of the above two steps are appended together
  into one phase oracle.
\end{itemize}

\subsection{ESOP Optimization to reduce T-count and T-depth}
\label{Pro:Minimize}

To minimize the number of $n \geq 4$ cubes in the ESOP, the proposal utilizes ESOP
minimization, such as EXORCISM-4~\cite{bib-exorcism}, to reduce the number of cubes in the
ESOP expression, thus reducing the T-count.

EXORCISM-4 is a heuristic method of minimizing the number of cubes in a given ESOP
expression. Extended discussion of its details is deferred to the cited paper.
Such a minimization not only reduces the number of cubes in the ESOP, it also increases
the distance between them i.e. reduces the number of shared variables, which enables more
parallelization when generating ESOP expressions. EXORCISM-4 also emphasizes

%% We modify it for our purposes by changing
%% the cost function \texttt{CountCubesInExactPseudoKro} that EXORCISM-4 uses for its heuristic, which uses a plain sum
%% of cubes to assess the cost of the arrangement that its heuristic encounters. We introduce a cost function
%% that instead uses a weighted sum of the cubes, where each cube is weighted by its T-count (i.e. a 2 literal cube
%% has weight 0) for $n \leq 3$ and $2^n$ for $n \geq 4$ (its T-count when synthesized according to the method in
%% Sec.~\ref{Pro:n4}). This penalizes terms of $n \geq 4$, while terms of $n < 3$ are free.


\subsection{Synthesize logic network and QC from $n \geq 4$ cubes}
\label{Pro:n4}

Because ESOP cubes can each be implemented as separate Controlled-Z gates, this
means that we can take each of the cubes in the ESOP expression and create a subject
graph with each of those cubes as separate outputs, such that it has as many
outputs as the ESOP expression has cubes. We can then construct a preliminary subject
graph in the following manner.

\begin{itemize}
\item Separate any cubes with $n \leq 3$ literals. Synthesize these directly from the 3-bit library.
\item For any cubes $n \geq 4$ literals, generate a cascade of AND nodes, with the first node connecting
  to two primary inputs, each node after that connecting to its output, and a different primary input,
  and the last node has its output directly to the final output.
\end{itemize}

This can then be fed into a logic network mapping tool such as abc~\cite{bib-mishchenko-lut} to
simplify Boolean logic network. Because this paper focuses on Toffoli gates and Controlled-Z gates,
we ask the logic network mapping tool to simplify our subject graph. Here, we do not specify
the simplification procedure: any such simplification procedure that optimizes for area such as
\texttt{mfs}~\cite{bib-mishchenko-mfs} or \texttt{lutpack}~\cite{bib-mishchenko-lutpack} will suffice.
We leave it to a topic of future research to expand on this degree of freedom. Meanwhile, we detail
which one we pick in our section on the implementation in~\secref{Exp}.

This Boolean logic network is processed just as in~\secref{Subsec:lhrs}. However, when it
needs to process a node directly connected to an output, it instead uses a phase gate
from the library we generated in the first step of~\secref{Subsec:ProMet}.

\begin{example}
  The steps of this example are depicted in~\figref{fig-ex-pro}.
  First we start off with the ESOP expression $\bar{x_{1}} \bar{x_{2}} x_{6} \bar{x_{5}} \bar{x_{4}} \oplus (x_{3} x_{6} \oplus x_{1} x_{5} \oplus \bar{x_{1}} \bar{x_{6}} \bar{x_{4}} \bar{x_{3}} \bar{x_{5}} \oplus \bar{x_{5}} \bar{x_{2}} \bar{x_{6}} x_{1} \bar{x_{4}} x_{3}$
  This is fed into an ESOP optimizer so we get
  $(x_{1} x_{5}) \oplus (\bar{x_{1}} \bar{x_{2}} x_{3} \bar{x_{4}}) \oplus (x_{1} \bar{x_{2}} \bar{x_{3}} x_{4} \bar{x_{5}}) \oplus x_{3}$.
  $x_1 x_5$ , and $x_3$ are generated directly as a Controlled-Z and a Z gate respectively, depicted
  in~\figref{fig-ex-qc}, in the $n \leq 3$ box.
  $(\bar{x_{1}} \bar{x_{2}} x_{3} \bar{x_{4}}) \oplus (x_{1} \bar{x_{2}} \bar{x_{3}} x_{4} \bar{x_{5}})$ are
  generated first as a subject graph, where each of the cubes $(\bar{x_{1}} \bar{x_{2}} x_{3} \bar{x_{4}})$
  and $(x_{1} \bar{x_{2}} \bar{x_{3}} x_{4} \bar{x_{5}})$ are
  depicted as a separate cascade of nodes, depicted in~\figref{fig-ex-network-og}.
  This subject graph is then fed to an area minimizer such as \texttt{abc}, resulting in~\figref{fig-ex-network},
  which has one less node due to sharing with the logic of $x_0$ and $x_1$.
  This subject graph is then processed such as in~\secref{Pro:n4}, resulting in the gates in the $n \geq 4$
  box in~\figref{fig-ex-qc}.
\end{example}

